\section{CoCa：对比学习之外再加一个生成头}

CLIP 已经很强，但它还是一个纯对比编码器。讲者接着引入 CoCa，思路很直接：如果只做对齐，模型会偏向“会认”；如果再加 captioning 生成目标，模型就必须“会说”，而会说通常意味着必须保留更多细粒度语义。\footnote{字幕 00:19:27--00:19:50 对 CoCa 的动机做了直接说明。}

\begin{figure}[H]
\centering
\includegraphics[width=0.95\textwidth]{slides-images/slide-042.jpg}
\caption{CoCa 在 CLIP 上加入生成式 caption loss}
\end{figure}

\begin{figure}[H]
\centering
\includegraphics[width=0.95\textwidth]{slides-images/slide-043.jpg}
\caption{CoCa 的结构：编码器负责对比学习，decoder 负责生成 caption}
\end{figure}

\begin{knowledgebox}{为什么要加 captioning}
对比学习擅长对齐图文，但生成目标会迫使模型保留更多区域级、属性级和关系级语义。CoCa 的思路是：既要对齐，也要会说。
\end{knowledgebox}

\begin{importantbox}{CoCa 的直觉}
把 captioning 当作一种更难的自监督约束，比只做 image-text matching 更能逼出 richer features。也就是说，它不是替代对比学习，而是在它上面多加一层语言压力。
\end{importantbox}

